% !TeX root = ../../Book1.tex
\section{交错型与外代数}
\label{b1:sec:B102}

有向面积不仅随边向量多线性变化，还应当在两条边相同的时候消失。交换两条边会反转方向，给某条边加上另一条边的倍数则不改变面积。行列式已经具有这些性质；外代数把它们从最高维推广到每一个维数，使一条有向线段、一个有向平行四边形以及高维平行多面体能够使用相同的运算规则。

\subsection{交错条件与行列式展开}
\label{b1:subsec:B10201}

\begin{definition}\label{b1:def:app-b-alternating-form}
设 \(k\geq1\)。若多线性型 \(\omega:V^k\rightarrow\mathbb R\) 在任意两个输入向量相同时取值为零，则称 \(\omega\) 为 \(V\) 上的 \(k\) 次\term[jiaocuo-xing]{交错型}[alternating form]。更一般地，对取值于向量空间的多线性映射，也用同一条件定义其交错性。零次交错型约定为实数。
\end{definition}

“交错型”一名可在黎景辉、白正简、周国晖的《高等线性代数学》第四章\cite{LiBaiZhou2014AdvancedLinear}中查到；该书第二、三章分别讨论张量积与张量代数，适合补充阅读本附录的代数构造。

交错条件蕴含交换两个输入会变号。固定其余输入，在所考察的两个位置同时代入 \(v+w\)，由多线性展开得到
\[
 0=\omega(\ldots,v,\ldots,w,\ldots)
   +\omega(\ldots,w,\ldots,v,\ldots).
\]
反过来，若交换两个输入会变号，则将二者取为同一个向量可得 \(2\omega=0\)，在实数域上这就给出交错性。因此任意排列 \(\sigma\) 满足
\[
 \omega(v_{\sigma(1)},\ldots,v_{\sigma(k)})
   =\operatorname{sgn}(\sigma)\omega(v_1,\ldots,v_k).
\]
若输入向量线性相关，把其中一个写成其他向量的线性组合，每项便含有重复输入，故值为零。尤其当 \(k>\dim V\) 时只有零交错型。

设 \(e_1,\ldots,e_n\) 是 \(V\) 的基，\(e^1,\ldots,e^n\) 为对偶基。对递增指标组 \(I=(i_1<\cdots<i_k)\)，定义
\[
 \epsilon^I(v_1,\ldots,v_k)
   =\det\bigl(e^{i_a}(v_b)\bigr)_{a,b=1}^k.
\]
这些都是交错型。把一般输入逐个按基展开，并把不同排列整理到递增次序，就得到
\begin{equation}\label{b1:eq:app-b-alternating-expansion}
 \omega(v_1,\ldots,v_k)
  =\sum_{i_1<\cdots<i_k}
       \omega(e_{i_1},\ldots,e_{i_k})
       \epsilon^I(v_1,\ldots,v_k).
\end{equation}
具体说，重复指标项为零；每一组互不相同的指标产生 \(k!\) 个项，其变号与坐标系数合起来恰是右侧的行列式。由于 \(\epsilon^I(e_{j_1},\ldots,e_{j_k})=\delta_{IJ}\)，这些型线性无关。于是 \(k\) 次交错型空间维数为 \(\binom nk\)，并由其在递增基向量组上的取值唯一决定。

\subsection{由张量代数取商}
\label{b1:subsec:B10202}

将各次张量积组成直和
\[
 T(V)=\bigoplus_{k=0}^{\infty}V^{\otimes k}.
\]
元素只含有限个非零分量。由连接张量因子定义乘法，零次的 \(1\in\mathbb R\) 为单位；结合律来自上一节的结合映射。这个代数称为\term[zhangliang-daishu]{张量代数}[tensor algebra]。它本身并不交换。

现在令 \(I\) 为 \(T(V)\) 中由所有 \(v\otimes v\) 生成的双边理想，也就是所有形如 \(a\otimes v\otimes v\otimes b\) 的元素的有限线性组合，其中 \(a,b\) 可分别属于任意次张量积。因为这些关系都有确定次数，按次数分组后仍是这类关系的和，所以 \(I\) 是各个 \(I\cap V^{\otimes k}\) 的直和。

\begin{definition}\label{b1:def:app-b-exterior-algebra}
上述商代数
\[
 \bigwedge V=T(V)/I
   =\bigoplus_{k=0}^{n}\bigwedge\nolimits^kV,
 \quad
 \bigwedge\nolimits^kV=
   V^{\otimes k}/\bigl(I\cap V^{\otimes k}\bigr)
\]
称为 \(V\) 的\term[waidaishu]{外代数}[exterior algebra]；其中 \(\bigwedge^kV\) 是它的 \(k\) 次分量。商中的乘法记为 \(\wedge\)，称为外积，纯张量的像记为 \(v_1\wedge\cdots\wedge v_k\)。约定 \(\bigwedge^0V=\mathbb R\)。
\end{definition}
\symidx{\(\bigwedge^kV\)：外代数的 \(k\) 次分量}

这里的外积与\subsecref{b1:subsec:170104}的外积是同一个构造；下面的基定理会说明它们的对应关系，故不重复登记“外积”的主索引。由商关系有 \(v\wedge v=0\)，再将 \(v+w\) 代入便得
\[
 v\wedge w=-w\wedge v.
\]
相邻两个因子可以交换变号；因此只要有两个因子相同，就可以将它们移到相邻位置并使整个乘积为零。商关系看似只处理相邻的重复向量，却已蕴含全部交错关系。

\begin{proposition}[外代数的基与交错泛性质]\label{b1:prop:app-b-exterior-basis}
若 \(e_1,\ldots,e_n\) 是 \(V\) 的基，则
\[
 e_I=e_{i_1}\wedge\cdots\wedge e_{i_k},
 \quad i_1<\cdots<i_k,
\]
组成 \(\bigwedge^kV\) 的基。特别地，
\[
 \dim\bigwedge\nolimits^kV=\binom nk,\quad
 \bigwedge\nolimits^1V=V,\quad
 \bigwedge\nolimits^kV=\{0\}\quad(k>n).
\]
每个交错多线性映射 \(B:V^k\rightarrow W\) 唯一对应线性映射
\(\widetilde B:\bigwedge^kV\rightarrow W\)，满足
\(\widetilde B(v_1\wedge\cdots\wedge v_k)=B(v_1,\ldots,v_k)\)。
\end{proposition}
\begin{proof}
把各向量按基展开，所得到的基向量乘积若有重复指标就为零，否则可以通过相邻交换整理成 \(\pm e_I\)。所以这些 \(e_I\) 张成 \(\bigwedge^kV\)。

为了证明无额外关系，对上面定义的 \(\epsilon^I\) 使用张量积的多线性泛性质，得到 \(V^{\otimes k}\) 上的线性泛函。凡是包含相邻的 \(v\otimes v\) 的生成关系，代入后两个输入相同，故它在 \(I\cap V^{\otimes k}\) 上为零。因此该泛函下降到 \(\bigwedge^kV\)，在 \(e_J\) 上的值为 \(\delta_{IJ}\)。逐一测试线性关系便得全部系数为零，证明线性无关。超过 \(n\) 个基向量必有重复，因而此时整个分量为零，也说明定义中的直和确实到 \(n\) 次为止。

最后对 \(B\) 作张量线性化。它同样在所有含相邻重复因子的关系上为零，所以下降为 \(\widetilde B\)。外积 \(v_1\wedge\cdots\wedge v_k\) 张成全空间，故下降后的映射唯一。反方向把任一线性 \(\widetilde B\) 与外积复合，所得映射具有多线性及交错性，因而两种操作互为逆。
\end{proof}

如果 \(v_j=\sum_i t_{ij}e_i\)，刚才的证明还给出坐标公式
\[
 v_1\wedge\cdots\wedge v_k
   =\sum_{|I|=k}\det(T_I)e_I.
\]
这正是\chapref{b1:ch:17}的子式构造。外代数的取商定义说明该公式为何能在任意基下成立；子式公式则提供实际计算的方法。两种表述各自解决一个问题，不需要在“抽象”与“计算”之间二选一。

外代数中形如 \(v_1\wedge\cdots\wedge v_k\) 的元素称为\term[kefenjie-waixiangliang]{可分解外向量}[decomposable multivector]；零元素也包括在内。并非所有外向量都可分解。对 \(\mathbb R^4\) 的标准基，令
\[
 \xi=e_1\wedge e_2+e_3\wedge e_4.
\]
则 \(\xi\wedge\xi=2e_1\wedge e_2\wedge e_3\wedge e_4\ne0\)，而任何 \(v\wedge w\) 与自身的外积都含重复因子，必为零。所以 \(\xi\) 不能表示成一个有向平行四边形。一般外向量更适合理解为这类元素的线性组合。

若 \(\xi\in\bigwedge^kV\)、\(\eta\in\bigwedge^\ell V\)，交换两个块需要交换 \(k\ell\) 对向量，故
\[
 \xi\wedge\eta=(-1)^{k\ell}\eta\wedge\xi.
\]
先对可分解元素证明，再利用双线性推广即可。这条规则没有说所有 \(\xi\wedge\xi\) 都为零；它只在 \(k\) 为奇数时强迫该结论。刚才的二次外向量例子正好展示偶数次数的差别。

\subsection{对偶配对与外积的归一化}
\label{b1:subsec:B10203}

对 \(V^*\) 也作同样的外代数构造。两类外空间的关系必须通过求值说明，不能因为都记作外积，就把向量和作用于向量的泛函混为一谈。

\begin{proposition}\label{b1:prop:app-b-form-pairing}
存在自然同构
\[
 \bigwedge\nolimits^kV^*
   \longrightarrow \bigl(\bigwedge\nolimits^kV\bigr)^*
\]
及其双线性配对，规定为
\[
 \left\langle
  \lambda_1\wedge\cdots\wedge\lambda_k,\,
  v_1\wedge\cdots\wedge v_k
 \right\rangle
 =\det\bigl(\lambda_i(v_j)\bigr)_{i,j=1}^k.
\]
经此同构，\(\bigwedge^kV^*\) 正是 \(V\) 上全部 \(k\) 次交错型的空间。
\end{proposition}
\begin{proof}
右侧对 \(\lambda_i\) 及 \(v_j\) 分别多线性、交错。先对一组变量使用外代数的泛性质，再对另一组使用，便得到良定义的双线性配对。在基
\[
 e^I=e^{i_1}\wedge\cdots\wedge e^{i_k},
 \quad e_J=e_{j_1}\wedge\cdots\wedge e_{j_k}
\]
上，配对值为 \(\delta_{IJ}\)。因此所定义的映射把 \(e^I\) 送到 \(e_I\) 的对偶基，故为同构。交错泛性质或公式\eqref{b1:eq:app-b-alternating-expansion}表明每个交错型都由这些基型的唯一线性组合给出。
\end{proof}
\symidx{\(\bigwedge^kV^*\)：\(V\) 上的 \(k\) 次交错型空间}

以后把 \(\omega(v_1,\ldots,v_k)\) 与
\(\langle\omega,v_1\wedge\cdots\wedge v_k\rangle\) 视为同一数。这里的尖括号是对偶配对；只有引入内积并作指定的识别后，才可进一步将它看作内积。

为了核对不同文献的系数，记 \(S_{k+\ell}\) 为排列群，并用
\(\operatorname{Sh}(k,\ell)\) 表示保持前 \(k\) 个位置及后 \(\ell\) 个位置各自顺序的排列。若 \(\alpha,\beta\) 的次数分别为 \(k,\ell\)，则由上述约定得到
\begin{equation}\label{b1:eq:app-b-wedge-shuffle}
 \begin{aligned}
 (\alpha\wedge\beta)(v_1,\ldots,v_{k+\ell})
  =\sum_{\sigma\in\operatorname{Sh}(k,\ell)}
     &\operatorname{sgn}(\sigma)\,
       \alpha(v_{\sigma(1)},\ldots,v_{\sigma(k)})\\
     &{}\cdot
       \beta(v_{\sigma(k+1)},\ldots,v_{\sigma(k+\ell)}).
 \end{aligned}
\end{equation}
证明只需先取 \(\alpha,\beta\) 为一形式的外积，将大行列式按前 \(k\) 行展开；每个被选中的列组对应一个保序排列，两个小行列式分别给出 \(\alpha,\beta\) 的值。一般情形再按外空间的基线性展开。等价地，可以对全部排列求和后除以 \(k!\ell!\)：同一保序排列对应的内部排列恰有 \(k!\ell!\) 个，内部变号分别被两个交错型吸收。

例如
\[
 (\lambda\wedge\mu)(v,w)
   =\lambda(v)\mu(w)-\lambda(w)\mu(v),
\]
没有额外的 \(\frac12\)。若把一般 \(r\) 次多线性型 \(T\) 的交错化定义为
\[
 \operatorname{Alt}T(v_1,\ldots,v_r)
  =\frac1{r!}\sum_{\sigma\in S_r}
       \operatorname{sgn}(\sigma)
       T(v_{\sigma(1)},\ldots,v_{\sigma(r)}),
\]
则本书的外积满足
\[
 \alpha\wedge\beta
   =\frac{(k+\ell)!}{k!\ell!}\,
      \operatorname{Alt}(\alpha\otimes\beta).
\]
这个因子来自两种不同的安排：\(\operatorname{Alt}\) 是平均，而外积要使基型与基外向量的配对恰为 \(1\)。以后积分微分形式时沿用后一约定。

类似地，把 \(\bigwedge^kV\) 识别为张量空间中的交错部分时，可送
\[
 v_1\wedge\cdots\wedge v_k
 \longmapsto
 \frac1{k!}\sum_{\sigma\in S_k}
   \operatorname{sgn}(\sigma)
   v_{\sigma(1)}\otimes\cdots\otimes v_{\sigma(k)}.
\]
在外代数商映射之后，这个表达又回到原来的外积，因此确为单射。这个识别中的 \(1/k!\) 与型的行列式求值公式并不冲突，但提醒我们：不能既使用张量子空间的平均定义，又未经换算套用另一种外积归一化。
